\documentclass[11pt]{article}
\usepackage{amsmath}
\usepackage[hidelinks]{hyperref}
\begin{document}
\title{Same late fields, different carried effects: when can a reply be packed?}
\author{}
\date{26 November 2026}
\maketitle

\paragraph{Audit.}
A late reply can carry an acknowledgment field and its data in
one byte instead of two separate byte attempts. That can keep
both the old report and the new reply within the same radio
budget. It cannot replace a distinct effect that an old executor
must consume before those data are ready. In the tested family two
plans can deliver the identical final decoded fields and still
give different physical clearance times. If the data are ready
early, a packed early reply can preserve both; otherwise a
separately authorized early path can repair the shared case.
These are authored R1/L2 source controls and exact same-information
direct matches, not an architecture or independently selected
non-radio result.

\paragraph{Protocol cut and authority.}
This experiment continues the site-readout issue in
\href{https://github.com/eternal-november/eos-research-programmes/pull/602}{PR 602}:
a mover may need a reliable indication of which site-side codebook
was used. Here the new application needs both labeled bits, so a
truthful common byte can carry both at once if the second bit is
available. Call the first bit and its current epoch a \emph{prelude}.
It is an authored application indication, not the transport
ACK mechanism or a mechanical completion receipt. The qualified
current endpoint interprets the type and epoch in the synthetic
payload. A stale prelude can arrive and spend bytes without being
authorized to satisfy the local condition.

In a prefix-sensitive mode the earlier old operator needs that
prelude by relative tick 2; a truthful reply at 4 is too late.
The test harness maps missing timely current prelude to an old
one-quantum hold before its next indexed move. The radio starts
at old simulation tick 8, so the prelude callback at relative 2
is at old tick 10. Physical prior completion comes from the
pinned warehouse trace with or without the previously demonstrated
idle step, not from the prelude itself. The late data are
ordinarily available no earlier than relative 4. The early-pack
control is admitted only when the data are independently
qualified available by 2. The two actors use fixed ordered
cuts and bounds in this fixture; neither the clock correspondence
nor the old authored actuation rule is established by the
SubT radio source. No other unmodeled old duty can be waived.

\paragraph{Actual shared resources and exact comparison.}
The carried prefix spends two bytes from P to Base. Its
mandatory old tail spends two more using the same endpoints
at relative tick 6. There is a fixed five-byte cap and all these
attempts occur within its unchanged one-second native rate window.
Each optional packet is one full application byte, even
though it has only a few semantic bits; this is not a claim
about field radio framing overhead. The prelude and data path
can independently alias the physical transmitter P or receiver
Base, or both, without changing logical recipients.
A separate independent native callback at tick 7 is required.
For a packed message, current prelude plus data actually
arrive in one callback. For split transport, the new context
consumes actual qualified prelude plus data, not the sender's
intention. A late packed callback cannot retroactively alter the
old source's earlier decision.

If the carried report must still go over P--Base, its exact
rate constraints at that step are
\[
2+U_P+2\leq5,\qquad 2+U_{Base}+2\leq5,
\]
where each $U$ counts actual optional debit on that endpoint
in the live window. They are separate inequalities. An
errored optional prelude spends its physical transmission
debit but not the receive debit; the experiment checks the
different receiver-only and sender-only overlap cases.
When both packets reach and both alias the same necessary
old endpoint, splitting loses the tail. Packing a delivered
reply can meet all the same deadlines where no earlier
consumption is necessary. However readiness for packing also
requires a common send time with $t_{data}\leq t_{send}
\leq t_{old\:cut}$ if that message is to perform the earlier
effect. When unavailable, a later reply is not a substitute
for the earlier use simply because both decoded fields agree.

\paragraph{Matched boundary.}
The actual old source clears at 18 or 19 in the selected
plain runs, and one hosted quantum changes those to 19 or 20.
Under the hard deadline 19, full reception and current
provenance, the exact choices are:
\begin{center}
\begin{tabular}{ll}
Declared fixture & Admitted whole plan \\
Shared ledgers, no earlier use, tight & Late packed \\
Disjoint ledgers, earlier use, tight & Split \\
Shared ledgers, earlier use, tight, no spare & None \\
Same, an authorized disjoint early path & Split via spare \\
Same, data actually ready at first slot & Early packed \\
Shared ledgers, earlier use, one tick slack & Late packed \\
\end{tabular}
\end{center}
In the last row the old completion is later but still within
the law: this is not equality of intermediate trajectories.
Likewise the spare is a real extra authority/resource premise,
not an advantage at an unchanged envelope. In the no-common-view
control one admitted source law needs the split plan while another
needs packing. An omniscient per-case selector cannot serve as
the earlier local implementation. A useful observation would
have to certify the consumed old effect or the relevant
aliases before the choice. With the spare allocated,
it is the only common choice across the two paired states.
A separately supplied lawful type receipt just before the
commitment lets the two states use their ordinary split or
pack choices without spending the spare. The same receipt a
tick later does not. This establishes the value of that
declared decision-specific observation, not the cost or
availability of obtaining it in a field deployment.

\paragraph{Evidence and limit.}
The native replay covers all 3,072 finite packet configurations,
including content, separate transmit/receive effects, age and
erasure. Twenty old physical profiles are measured and seven
same-process staged calls continue the radio execution after
running the selected old physical branch from the actual
prelude callback. 40,960 full interpreter/direct comparisons
agree, with four scoped Lean statements. The comparable direct
model gets the same packet choices and old effects, so the
evidence is an exact boundary relative to the specified
interfaces, not an improvement over an exact scheduler that
already makes those same distinctions. A stricter setting
could replace this old timing rule with a measured actuator
law or introduce a deadline-sensitive placement suitable for
independent testing. Neither is established by these runs.
\end{document}
